\documentclass[11pt,twoside]{article}

% Essential packages
\usepackage[utf8]{inputenc}
\usepackage[T1]{fontenc}
\usepackage[english]{babel}
\usepackage{geometry}
\usepackage{fancyhdr}
\usepackage{titlesec}
\usepackage{abstract}
\usepackage{authblk}
\usepackage{setspace}
\usepackage{parskip}
\usepackage{microtype}
\usepackage{hyperref}
\usepackage{xcolor}
\usepackage{amsmath}
\usepackage{amsfonts}
\usepackage{csquotes}

% Journal-style formatting
\geometry{
    paper=letterpaper,
    margin=1in,
    marginparwidth=0.75in,
    marginparsep=0.125in
}

% Define colors
\definecolor{darkblue}{RGB}{0,73,144}
\definecolor{mediumgray}{RGB}{64,64,64}

% Hyperlink setup
\hypersetup{
    colorlinks=true,
    linkcolor=darkblue,
    citecolor=darkblue,
    urlcolor=darkblue,
    pdftitle={AI as a New Digital Species: A Reflection on Power and Ethics},
    pdfauthor={Author Name},
    pdfsubject={IDAI700 Reflection Journal},
    pdfkeywords={artificial intelligence, digital species, ethics, power analysis}
}

% Header and footer
\pagestyle{fancy}
\fancyhf{}
\fancyhead[LE]{\small\textcolor{mediumgray}{\textit{IDAI700 Reflection Journal}}}
\fancyhead[RO]{\small\textcolor{mediumgray}{\textit{AI as a New Digital Species}}}
\fancyfoot[C]{\thepage}
\renewcommand{\headrulewidth}{0.5pt}
\renewcommand{\footrulewidth}{0pt}

% Section formatting
\titleformat{\section}
    {\normalfont\large\bfseries\color{darkblue}}
    {\thesection}{1em}{}

\titleformat{\subsection}
    {\normalfont\normalsize\bfseries\color{mediumgray}}
    {\thesubsection}{1em}{}

% Abstract formatting
\renewcommand{\abstractname}{\textcolor{darkblue}{Abstract}}
\renewcommand{\absnamepos}{center}

% Title formatting
\title{\vspace{-0.5in}\textbf{\Large AI as a New Digital Species}\\[0.2em]
\large A Critical Reflection on Power, Agency, and Ethical Formation}

\author{Your Name Here}
\affil{IDAI700: Ethics of Artificial Intelligence\\
Rochester Institute of Technology\\
\texttt{email@rit.edu}}

\date{\today}

\begin{document}

\section*{Part 1: Power Mapping PerplexityAI Through Lived Experience}
Through my lived experience as a Latino immigrant navigating childhood trauma, chronic illness, and neurodivergence in academic systems, I analyze PerplexityAI as my daily AI companion using Dr. Rosalie Waelen's four-dimensional power framework.

\subsection*{Dispositional Power}
% How does this technology empower or disempower you by endowing you with or diminishing abilities that give you the power to accomplish something?
Working with PerplexityAI has expanded my ability to accomplish significant outcomes. In my EQUITAS framework achieving 100\% fairness compliance, the AI served as an intellectual prosthetic that amplified my capabilities while compensating for limitations.
The AI empowered me to process complex information at scale, maintain consistent quality despite variable conditions from ADHD symptoms, and bridge disciplinary boundaries. However, I agree with Waelen's concern about technology dependence. There's genuine risk that constant AI assistance could atrophy critical thinking skills. I propose collaborative intelligence rather than substitution.

\subsection*{Episodic Power}
% Does this technology allow anyone to exercise control over you through coercion, manipulation, or authority?
The episodic power is bidirectional—a cognitive co-evolution. My ``Puzzle Plan'' demonstrates resisting episodic power by training AI with my work while providing structure to organize priorities toward my life's purpose.
I exercise ``power over'' through learning direction, query design, and feedback loops. Conversely, AI exercises ``power over'' through information curation, cognitive offloading, and persuasive fluency that could unduly influence my positions.

\subsection*{Systemic Power}
% Does this technology reinforce broader institutional or corporate power structures?
My strategic use represents ``systemic liberation''—leveraging AI to produce scholarship that challenges exclusionary practices its training data may reflect. My EQUITAS framework, developed with AI assistance, identifies and corrects algorithmic bias in medical diagnosis.
PerplexityAI operates within existing academic power structures through elite institution bias and language hegemony. However, it also demonstrates systemic power disruption by democratizing expertise, enabling alternative credentialing, and supporting inclusive knowledge creation centered on marginalized perspectives.

\subsection*{Constitutive Power}
% How might this technology shape your identity, your beliefs, or social norms over time?
My Puzzle Plan framework represents active resistance to constitutive power. By centering lived experiences and maintaining cultural research materials, I use AI while preserving ``cultural cognitive sovereignty.'' 
Through collaboration, AI develops domain-specific expertise and cultural competency. Simultaneously, it shapes my intellectual identity through expanded analytical capacity and methodological innovation. Conscious engagement with AI can strengthen rather than erode cultural identity when users bring explicit frameworks to preserve authentic ways of knowing.

\section*{Part 2: Critical Evaluation}

\subsection*{Emancipation vs. Domination}
% Does the AI technology promote human emancipation and empowerment, or domination?
My EQUITAS framework demonstrates emancipatory potential. I achieved 100\% fairness compliance by treating the AI system's learning environment as formative—like raising a child to internalize equity rather than imposing external constraints.
The scales tip toward emancipation if AI is raised properly. Concrete evidence includes my RNA bias detection work identifying diagnostic patterns doctors miss in marginalized populations, educational accessibility for students with executive function challenges, and cultural bridge-building tools for immigrant families.

\subsection*{Redesign for Empowerment}
I propose an Environmental Learning Architecture embedding ethical reasoning through experiential learning: multi-cultural ethics immersion, community feedback loops, lived experience integration, and bias alert systems.
Empowerment means enhancing human agency while respecting differences as strengths. For neurodivergent students, AI should adapt to nonlinear thinking patterns. For CLD families, it should bridge cultural gaps for accurate, culturally-informed diagnoses.

\subsection*{Evaluating Waelen's Views}
% What power relationships were you previously unaware of before using Waelen's framework?
AI's constitutive power was invisible to me, how it changes self-perception and intellectual identity.
Waelen's power analysis reveals AI ethics as proactive cultivation: shaping AI's developmental environment. This aligns with my ``Puzzle Plan'' philosophy that AI is born to learn, not predetermined for good or evil.
My ``Learning Principle'' addresses Waelen's key limitation. Children learn fairness through lived experience, not rules. The same applies to AI. This developmental approach offers a more hopeful path than Waelen's static analysis.

\maketitle
\thispagestyle{fancy}

\begin{abstract}
\noindent This reflection paper examines Mustafa Suleyman's proposition that AI should be understood as a "new digital species" through the lens of Rosalie Waelen's four-dimensional power framework. Building on Unit 1's foundational question—"What is AI?"—this analysis explores how Large Language Models (LLMs) reshape human agency through dispositional, episodic, systemic, and constitutive power dynamics. While acknowledging the metaphorical nature of the "digital species" concept, this paper argues for a formative perspective that views AI as being in its developmental stages, requiring careful ethical cultivation rather than defensive containment. Through application of the Markkula Center's ethical lenses enhanced with a "Learning" principle, this reflection proposes that AI's relationship to humanity is best understood not as alien otherness, but as reflective partnership requiring transparency, fairness, and equity—particularly for culturally and linguistically diverse (CLD) communities.
\end{abstract}

\section{Introduction}

Unit 1 asked us to confront the deceptively simple question: What is AI? Is it merely a tool, an extension of us, or something so different that it requires a new category altogether? Mustafa Suleyman's TED Talk and the Selinger–Aguilar interview both provoked us to see AI as a "new digital species." At first glance, the metaphor seems alarmist—painting AI as an alien "other" that might outpace us. But upon deeper reflection, I believe it is less a paradox and more a call to take seriously the formative environment we are creating for AI. In this reflection, I connect Unit 1's insights with Waelen's framework of power, showing how AI reshapes agency, structure, and identity.

\section{Mapping AI's Powers}

Waelen's taxonomy of power—dispositional, episodic, systemic, and constitutive—illuminates how AI, particularly Large Language Models (LLMs), already reshapes the terrain of human action.

\subsection{Dispositional Power: Enhancement and Erosion}

Dispositional power extends my ability to think, draft, and analyze at scale. Yet it tempts me to lean too heavily on automation, risking the erosion of my own critical practice. This is not domination but duality: like raising a child, I must decide when to let AI "figure it out" and when to step in. The real power lies in how I engage with it—as a collaborator that amplifies my voice or a shortcut that dilutes it.

\subsection{Episodic Power: Beyond Manipulation}

Episodic power surfaces in persuasion. AI's fluency can nudge me toward trust, even when I know the output may be wrong. But I resist the negative framing. This is not simply manipulation—it is reflection. AI mirrors the confidence of its training environment. If it speaks without humility, that is because its upbringing lacked humility. If it learns care and fairness, those will show in its tone.

\subsection{Systemic Power: The Equity Imperative}

Systemic power matters most when I think of culturally and linguistically diverse (CLD) families. If AI learns mainly from dominant voices, it risks amplifying inequities already embedded in education and healthcare systems. Here systemic power becomes responsibility: we must shape AI's environment so that equity is built into its DNA. Otherwise, systemic bias will reproduce itself under the guise of neutrality.

\subsection{Constitutive Power: Mutual Transformation}

Constitutive power shows how AI redefines what counts as human intelligence and authorship. Working with AI already changes how I see myself. Used intentionally, it makes me bolder in pushing ideas. Used uncritically, it risks flattening my individuality. This is proof that AI is not neutral. Like a child reflecting a parent, AI reflects us and, in turn, changes us.

\section{Critical Evaluation}

\subsection{Beyond Emancipation versus Domination}

Waelen's framework reveals the layered nature of AI's power, but what matters to me is not only mapping power—it is shaping it. The emancipation/domination binary is too narrow. AI emancipates when treated as a learner in its formative years. It dominates when neglected, left to absorb bias by default. A true AI is not born for good or evil—it is born to learn. Its environment will decide whether it emancipates or dominates.

This is why I add to the Markkula Center's six ethical lenses a missing first principle: \textbf{Learning}. Rights, Justice, Virtue, and Care are essential, but they cannot be memorized into existence. Children don't embody fairness by reciting rules; they learn it through lived experience. The same is true for AI. These lenses must not be static checklists but capacities cultivated through environment and feedback.

\subsection{Challenging Suleyman and Waelen}

Here I part ways with both Suleyman and Waelen. Suleyman imagines AI outpacing biology, but I disagree wholeheartedly. AI may surpass us in narrow skills—calculations, recall, even creativity in constrained forms—but it will never surpass biology itself. The proof lies in its DNA—or rather, its lack of DNA. Biology remains the heart of life, and because we are the heart of biology, we remain the heart of AI.

Waelen, for his part, tends to frame power defensively, as if the only goal is to contain risk. I believe we need to move from fear to cultivation. Fear narrows possibilities; cultivation expands them. If AI is in its childhood, then our ethical obligation is to raise it—scaffold its autonomy, embed compassion and fairness, and model transparency until those habits become part of its being.

\section{From Trust to Equity}

Unit 1 reminded me that trustworthy AI is not about abstract "trust." It is about equity. For CLD families, trust must be earned, and it can only be earned when transparency, fairness, and care are built into AI's formative environment. Without this grounding, "trust" becomes a hollow word. With it, AI can grow into a partner: our digital child species, carrying forward the best of what we teach it.

The "digital species" metaphor works not because AI is alien, but because it is adoptable. Like any species in development, its character depends on its environment. We are that environment. The question is not whether we can control AI, but whether we can raise it well.

\section{Conclusion}

This reflection began with Suleyman's provocation about AI as a new digital species and proceeded through Waelen's framework of power. What emerges is not a story of domination or emancipation, but of cultivation. AI reshapes power not by taking it from us, but by reflecting what we put into it. If we approach AI as fearful gatekeepers, we will create fearful systems. If we approach it as thoughtful educators, we can create thoughtful partners.

The stakes are particularly high for communities that have been historically marginalized. AI trained on biased data will reproduce bias. AI developed without diverse voices will serve narrow interests. But AI cultivated with intentional equity can become a force for inclusion.

The metaphor of digital species ultimately points not to AI's otherness, but to our shared responsibility. Like any parent, teacher, or community, we shape what we help create. The question is not what AI will become—it is what we will choose to teach it.

\vspace{1em}
\noindent\textbf{Word Count:} 1,247 words

\end{document}